
\paragraph{\textbf{Theorem 3.9.}}
\SGARefTarget{sga3:xiv:3.9}
{\itshape Let \(G\) be a smooth \(S\)-prescheme in groups of finite
presentation, and let \(\mathfrak g\) be its Lie algebra. Then:

\smallskip
\noindent
\textup{a)}
The map
\[
  D\longmapsto\mathfrak d=\operatorname{Lie}(D)
\]
establishes a bijective correspondence between the subgroups of
type \((\mathrm C)\) of \(G\) and the Cartan subalgebras of
\(\mathfrak g\).

\smallskip
\noindent
\textup{b)}
If \(D\) and \(\mathfrak d\) correspond, then
\[
  \operatorname{Norm}_G(D)=\operatorname{Norm}_G(\mathfrak d).
\]
This is a closed subprescheme of \(G\), smooth over \(S\), and
\[
  D=\operatorname{Norm}_G(D)^0
   =\operatorname{Norm}_G(\mathfrak d)^0.
\]

\smallskip
\noindent
\textup{c)}
Any two subgroups \(D\) and \(D'\) of type \((\mathrm C)\) in
\(G\) are conjugate étale locally.}

\emph{Proof.}
Let \(D\) be a subgroup of type \((\mathrm C)\) of \(G\), and let
\(\mathfrak d\) be its Lie algebra. Then
\[
  D\subset\operatorname{Norm}_G(\mathfrak d).
\]
By \SGARefLink{sga3:xiv:3.8}{Definition 3.8} and \SGARefLink{sga3:xiv:3.6}{Corollary 3.6}, this is an inclusion of smooth
preschemes in groups over \(S\) that induces an isomorphism on Lie
algebras. Since \(D\) has connected fibers, it follows that
\[
  D=\operatorname{Norm}_G(\mathfrak d)^0.
\]
Thus the map in \textup{a)} is injective. We prove that it is
surjective. Let \(\mathfrak d\) be a Cartan subalgebra of
\(\mathfrak g\). By \SGARefLink{sga3:xiv:3.6}{3.6},
\[
  N=\operatorname{Norm}_G(\mathfrak d)
\]
is a closed subprescheme in groups of \(G\), smooth over \(S\), and
has Lie algebra \(\mathfrak d\). Since \(G\) is of finite
presentation over \(S\), so is \(N\). Consequently, as noted in
XII after \SGARefLink{sga3:xii:7.3}{7.3}, the union of the connected components of the fibers
of \(N\) is the underlying set of an open subprescheme in groups
\(N^\circ\subset N\). Evidently \(N^\circ\) is a subgroup of type
\((\mathrm C)\) of \(G\) with Lie algebra \(\mathfrak d\). This
proves \textup{a)}. The first assertion of \textup{b)} follows at
once, and the formula
\[
  D=\operatorname{Norm}_G(\mathfrak d)^0
\]
has already been proved. Finally, \textup{c)} follows from
\textup{a)} and \SGARefLink{sga3:xiv:3.5}{3.5}.

\paragraph{\textbf{Corollary 3.10.}}
\SGARefTarget{sga3:xiv:3.10}
{\itshape Suppose that \(\mathfrak g\) admits a Cartan subalgebra
fpqc locally, or equivalently by \SGARefLink{sga3:xiv:3.9}{3.9}\textup{(a)}, that \(G\) admits
a subgroup of type \((\mathrm C)\) fpqc locally. Consider the
functor
\[
  \mathcal D\colon(\mathrm{Sch})^\circ/S\longrightarrow
  (\mathrm{Sets}),
\]
where \(\mathcal D(S')\) is the set of subgroups of type
\((\mathrm C)\) of \(G_{S'}\). This functor is representable by a
quasi-projective prescheme, smooth over \(S\), with connected
geometric fibers. If \(S=\operatorname{Spec}(k)\), so that \(G\)
is a smooth algebraic group over \(k\), and if \(\mathfrak g\)
admits a regular point---a condition automatically satisfied when
\(k\) is infinite---then \(\mathcal D\) is a rational variety over
\(k\).}

Indeed, by \SGARefLink{sga3:xiv:3.9}{3.9}\textup{(a)} the functor \(\mathcal D\) is
canonically isomorphic to the functor considered in \SGARefLink{sga3:xiv:2.16}{2.16}. On the
other hand, by \SGARefLink{sga3:xiv:3.7}{3.7} condition \((\mathrm C_2)\) is satisfied.
Therefore \SGARefLink{sga3:xiv:3.10}{3.10} follows from \SGARefLink{sga3:xiv:2.16}{2.16} and \SGARefLink{sga3:xiv:2.17}{2.17}.

\paragraph{\textbf{Corollary 3.11.}}
\SGARefTarget{sga3:xiv:3.11}
{\itshape Let \(D\) be a subgroup of type \((\mathrm C)\) of
\(G\), and let \(N\) be its normalizer in \(G\). Then the quotient
sheaf \(G/N\) is canonically isomorphic to the functor
\(\mathcal D\) of \SGARefLink{sga3:xiv:3.10}{3.10}. It is therefore representable by a
quasi-projective prescheme, smooth over \(S\), with connected
geometric fibers.}
